\documentclass[../../../include/open-logic-section]{subfiles}

\begin{document}

\olfileid{sth}{story}{predicative}	

\olsection{స్వవర్గానవలంబిత, స్వవర్గావలంబిత నిర్వచనాలు}

రసెల్ సమితి~$R$ను
$\Setabs{x}{x \notin
x}$తో నిర్వచించారు. మరింత
వివరంగా చెప్పాలంటే, $R$
అనేది తమలో తాము మూలకాలుగా
లేని సమితులన్నింటినీ,
వాటినే కలిగిన సమితి.
\sourcecorrection{OLTESTPRED-001}{మూల విస్తృత వివరణలో ‘తమలో తాము మూలకాలుగా లేని సమితులు కానివి’ అని ద్వి-నిరాకరణతో రాసింది; అది $x \notin x$ నిర్వచనానికి విరుద్ధం. నిర్వచనానికి సరిపోయే ‘తమలో తాము మూలకాలుగా లేని సమితులు’గా సరిచేశాం.}
అందువల్ల $R$ను నిర్వచించేటప్పుడు,
అది ఉంటే~$R$నూ కలిగే
ప్రవేశంపై మనం పరిమాణకం
వేస్తున్నాం.

ఇది \emph{స్వవర్గావలంబిత
(impredicative)} నిర్వచనం.
సాధారణంగా, నిర్వచించబడే
వస్తువునే కలిగిన ప్రవేశంపై
పరిమాణకం వేసినప్పుడే
ఒక నిర్వచనం స్వవర్గావలంబితమని
చెప్పవచ్చు.
	
ఈ వైరుధ్యాల తరువాత వైట్‌హెడ్,
రసెల్, పువాంకారే, వైల్
ఇలాంటి స్వవర్గావలంబిత
నిర్వచనాలను ``దుష్టవలయాలు''గా
తిరస్కరించారు:
\begin{quote}
	నివారించాల్సిన వైరుధ్యాలను
	విశ్లేషిస్తే, అవన్నీ ఒక రకమైన
	దుష్టవలయం నుంచి వచ్చాయని
	తెలుస్తుంది. ఒక వస్తు సముదాయంలో,
	ఆ సముదాయం మొత్తాన్ని ఆధారంగా
	తీసుకుంటేనే నిర్వచించగల
	సభ్యులు ఉండవచ్చని
	అనుకోవడమే ఈ దుష్టవలయాలకు
	మూలం[\ldots. \textparagraph]

	అనుచిత సమగ్రతలను నివారించే
	సూత్రాన్ని ఇలా చెప్పవచ్చు:
	`ఒక సముదాయంలోని \emph{అన్నిటినీ}
	ఆధారంగా తీసుకునేది ఆ
	సముదాయంలో ఒకటిగా ఉండకూడదు';
	లేదా తిరిగి చెప్పాలంటే:
	`ఒక సముదాయానికి సంపూర్ణం
	ఉందనుకుంటే, ఆ సంపూర్ణాన్ని
	ఆధారంగా తీసుకుంటేనే
	నిర్వచించగల సభ్యులు దానిలో
	ఉంటే, ఆ సముదాయానికి
	సంపూర్ణం లేదు.' దీనిని
	`దుష్టవలయ సూత్రం' అంటాం;
	ఎందుకంటే అనుచిత సమగ్రతలను
	అనుకోవడం వల్ల వచ్చే
	దుష్టవలయాలను ఇది
	నివారిస్తుంది.
	\citep[p.~37]{WhiteheadRussell1910}
\end{quote}
వారి మాదిరిగా \emph{దుష్టవలయ
సూత్రాన్ని అంగీకరిస్తే},
\olref[sth][story][rus]{sec}లోని
వినాశకర నిర్బంధరహిత
ధర్మసంగ్రహ పథకం స్థానంలో
ఇలాంటిదాన్ని పెట్టవచ్చు:
\sourcecorrection{OLTESTPRED-002}{మూల అనుసంధాన వాక్యం స్వవర్గావలంబిత నిర్వచనాలను తిరస్కరించిన రచయితలను అనుసరిస్తూనే దుష్టవలయ సూత్రాన్ని ‘తిరస్కరిస్తే’ అంది. వారిని అనుసరించడమంటే ఆ సూత్రాన్ని అంగీకరించడమే; తదుపరి రక-స్థాయిల నిర్మాణానికి అనుగుణంగా సరిచేశాం.}

\begin{defish}
\emph{స్వవర్గానవలంబిత ధర్మసంగ్రహం.}
సమితులపై మాత్రమే పరిమాణకం
వేసే ప్రతి సూత్రం $\phi$కు:
సమితి$^\prime$ $\Setabs{x}{\phi(x)}$
ఉంటుంది.
\end{defish}

సమితులు$^{\prime}$ సాధారణ
సమితులు కానంతవరకు
వైరుధ్యం రాదు.

దురదృష్టవశాత్తూ,
స్వవర్గానవలంబిత ధర్మసంగ్రహం
అంత \emph{సమగ్రం} కాదు.
అది కొత్త వస్తువులను,
సమితులు$^\prime$ను,
పరిచయం చేస్తుంది. కాబట్టి
సమితులు$^\prime$పై
పరిమాణకం వేసే సూత్రాలను
కూడా పరిశీలించాలి. అవి
ఎప్పుడూ సమితి$^\prime$నే
ఇస్తే, తమలో తాము మూలకాలుగా
లేని సమితులు$^\prime$
అన్నింటి సమితి$^\prime$ను
పరిశీలించగానే రసెల్
వైరుధ్యం మళ్లీ వస్తుంది.
అందువల్ల అదే భావాన్ని
కొనసాగిస్తే, సమితులు$^\prime$పై
పరిమాణకం వేసే సూత్రం
సంబంధిత సమితి$^{\prime\prime}$ను
ఇస్తుందని చెప్పాలి. ఆ తరువాత
సమితులు$^{\prime\prime\prime}$,
సమితులు$^{\prime\prime\prime\prime}$,
వగైరా అవసరం. ప్రైమ్ గుర్తులు
పెరిగిపోకుండా, వీటిని
సమితులు$_0$, సమితులు$_1$,
సమితులు$_2$, సమితులు$_3$,
సమితులు$_4$,\ldots గా
భావించడం సులభం. ఇది
(సరళ) రకాల సిద్ధాంతానికి
దారి చూపుతుంది.

అలాంటి సిద్ధాంతంపై కొన్ని
స్పష్టమైన అభ్యంతరాలున్నాయి
(అయితే అవి \emph{నిర్ణాయక}
అభ్యంతరాలేనని స్పష్టం కాదు).
సంక్షిప్తంగా: ఫలిత సిద్ధాంతాన్ని
వాడడం క్లిష్టం; అనేక రకాల
వస్తువులను అనవసరంగా
ప్రతిపాదిస్తుంది; చివరికి
స్వవర్గావలంబిత నిర్వచనాలు
నిజంగా \emph{అంత చెడ్డవేనా}
అన్నదీ స్పష్టం కాదు.
	
చివరి అంశాన్ని చూపడానికి
\citeauthor{Ramsey1925} చేసిన
ఈ వ్యాఖ్యను పరిశీలించండి:
\begin{quote}
	ఒక సమూహంలోని అత్యంత
	పొడవైన వ్యక్తి అని
	ఒకరిని పేర్కొనవచ్చు;
	ఆయన స్వయంగా సభ్యుడైన
	మొత్తం సమూహాన్ని
	ఆధారంగా తీసుకొని
	ఆయనను గుర్తించినా,
	అందులో దుష్టవలయం
	ఏమీ లేదు. \citep{Ramsey1925}
\end{quote}
రామ్సే ఉద్దేశం: ``సమూహంలోని
అత్యంత పొడవైన వ్యక్తి''
\emph{స్వవర్గావలంబిత}
నిర్వచనమే; అయినా అది
స్పష్టంగా నిర్దోషమైనది.

ఈ సందర్భంలో అత్యంత
పొడవైన వ్యక్తిని
\emph{స్వవర్గానవలంబిత}
పద్ధతిలోనూ గుర్తించవచ్చని
ఎవరైనా బదులివ్వవచ్చు.
ఉదాహరణకు, ఆ వ్యక్తిని
నేరుగా వేలితో చూపవచ్చు.
అప్పుడు స్వవర్గావలంబిత
నిర్వచనాలపై అభ్యంతరం,
\emph{కేవలం} అలాంటి
పద్ధతిలోనే గుర్తించగల
వస్తువులకు పరిమితమవాలి.
అయినా అలాంటి ``అత్యవసర
స్వవర్గావలంబన'' ఎందుకు
చెడ్డదో మరింత వివరణ
కావాలి.\footnote{మరిన్ని
వివరాలకు \citet{Linnebo2010}
చూడండి.}

వస్తువును \emph{సృష్టించడానికి}
ఒక వంటకం లాంటిది నిర్వచనం
ఇవ్వాలని అనుకుంటే మాత్రం,
స్వవర్గావలంబిత నిర్వచనాలు
నిజంగానే తీవ్రమైన సమస్య.
అలాంటి నిర్వచనంలో
దుష్టవలయంలో చిక్కుకుంటాం:
స్వవర్గావలంబితంగా పేర్కొన్న
వస్తువును సృష్టించాలంటే,
ఆ వస్తువుతో సహా \emph{అన్ని}
వస్తువులనూ \emph{మొదట}
సృష్టించాలి; ఎందుకంటే
దానిని అన్ని వస్తువుల
ఆధారంగా పేర్కొన్నారు.
అంటే ఆ వస్తువును
సృష్టించకముందే దానిని
సృష్టించాల్సి వస్తుంది.
కానీ సమితులను మనం
\emph{సృష్టించే} వస్తువులుగా
భావించినప్పుడే ఇది
``అత్యవసర స్వవర్గావలంబన''తో
నిర్వచించిన సమితులపై
తీవ్రమైన అభ్యంతరం.
బహుశా మనం అలా భావించం.

అందువల్ల---మంచికో చెడుకో---
ప్రచారంలోకి వచ్చిన పద్ధతి
స్వవర్గానవలంబన లేదా
స్వవర్గావలంబన విషయంలో
కఠిన వైఖరి తీసుకోదు.
దాని బదులు, ఇప్పుడు
సంచిత-దశలవారీ పద్ధతి
అని భావించే దానిని వాడుతుంది.
చివరకు ఇది మన సమితులను
``దశలుగా'' అమర్చడానికి
వీలు కల్పిస్తుంది---
స్వవర్గానవలంబిత పద్ధతి
వస్తువులను సమితులు$_0$,
సమితులు$_1$, సమితులు$_2$,
\ldots గా అమర్చినట్లుగా
\emph{కొంతవరకు}---కానీ
వాటి మధ్య రకపరమైన
భేదాన్ని ప్రతిపాదించం.

\end{document}
